\scanpage{38}
\begin{center}
UNIVERSITÉ DE PARIS

\figure{\begin{tikzpicture}\draw (-.4,0)--(-.06,0);\fill (0,0) circle (.7pt);\draw (.06,0)--(.4,0);\end{tikzpicture}}

ÉCOLE NORMALE SUPÉRIEURE

\figure{\begin{tikzpicture}\draw (-.4,0)--(-.06,0);\fill (0,0) circle (.7pt);\draw (.06,0)--(.4,0);\end{tikzpicture}}

45, RUE D'ULM

ODÉON 06$\cdot$45

\figure{\begin{tikzpicture}\draw (-.4,0)--(-.06,0);\fill (0,0) circle (.7pt);\draw (.06,0)--(.4,0);\end{tikzpicture}}
\end{center}

\uuline{Dém}. $G/S$ opère sans point fixe sur $E/S$ : si
\[
gx=s_0x,\quad g=s_0s\text{ avec }sx=x\quad\Rightarrow\quad g\in S.
\]
Si $g\notin S$, $g$ transforme tout sommet de $E$ en un sommet distinct, donc tout simplexe en un simplexe disjoint. $E/G=E/S\big/G/S$.

---

\uline{Le Théorème 1} \deleted{\unsure{en}} résulte de ce lemme et de ce qui a été vu précédemment.

\deleted{$\unsure{p_*\pi_1(B)|p_*\pi_1(E)}$}

\mspage{37}

